% 16_generation_examples.tex -- appendix: prompt + setup examples
\section{Prompt and Generation Examples}
\label{app:examples}

This appendix shows representative prompts from each of the four
content domains used in \autoref{sec:experiments}, together with the
exact ChainMark configuration and post-attack pipeline that produced
the headline numbers in \autoref{tab:headline-aggregate} and
\autoref{tab:headline-percell}. All three models receive the same
prompt verbatim with no system prompt prepended; decoding parameters
($T{=}0.7$, top-$p{=}1$, $n{=}200$ tokens, deterministic
prompt-order seed) are also held fixed across models and methods.

\paragraph{Prompt examples per domain.}
\begin{itemize}[nosep, leftmargin=*]
  \item \textbf{Code} (HumanEval problem stems, $164$ problems
    available; we use the first $100$):
    \begin{quote}\small\ttfamily
    from typing import List\\
    def has\_close\_elements(numbers: List[float], threshold: float) -> bool:\\
    \quad """Check if in given list of numbers, are any two numbers closer to each other than given threshold."""
    \end{quote}
  \item \textbf{Factual} (short closed-form knowledge prompts):
    \begin{quote}\small\ttfamily
    The capital of France is
    \end{quote}
  \item \textbf{Wiki} (open-ended ``Explain $X$\ldots'' prompts over
    a curated $176$-entry concept list):
    \begin{quote}\small\ttfamily
    Explain Donald Trump in a comprehensive way.
    \end{quote}
  \item \textbf{Writing} (creative-completion prompts):
    \begin{quote}\small\ttfamily
    Write a short story that begins: The colony ship arrived three centuries late, and someone was already waiting.
    \end{quote}
\end{itemize}

\paragraph{ChainMark generation configuration (used in
\autoref{tab:headline-aggregate}).}
At each generation step on the prompt above, the model produces a
distribution over its native vocabulary; ChainMark then masks logits
at gated positions (gate density $\rho{=}0.5$, high-entropy gate
$G_{H_{\mathrm{high}}}$) to keep only token IDs whose state
$\sigma_\kappa(t) = \mathrm{SHA\text{-}256}(\kappa\,\Vert\,t) \bmod
S$ equals $(s+1) \bmod S$, where $s$ is the previous token's state.
At ungated positions the model samples at $T{=}0.7$ from the original
distribution. The secret key $\kappa$ is fixed for the full
campaign; in deployment a regulator-side audit re-derives the same
state sequence from the same key with the model's tokenizer.

\paragraph{Attack-pipeline example.} Given the watermarked text $x$
(say, $200$ tokens of a wiki-domain answer), the three attack
conditions in \autoref{tab:headline-aggregate} apply respectively:
\begin{itemize}[nosep, leftmargin=*]
  \item \textbf{Clean.} Detect on $x$ as-is.
  \item \textbf{Random substitution at $\delta_{\mathrm{eff}}{=}0.20$.}
    Replace a uniformly random $20\%$ of tokens with a uniformly
    sampled vocabulary token, then detect on the corrupted text.
  \item \textbf{ZH back-translation.} Pass $x$ through
    NLLB-200~\citep{nllbteam2022nllb}
    English\,$\rightarrow$\,Chinese\,$\rightarrow$\,English; record the
    empirical token-edit-distance $\delta_{\mathrm{eff}}$
    ($\approx 0.81$ for ZH on average,
    \autoref{tab:translation-pivots}); then detect on the
    round-tripped text.
\end{itemize}
The detector runs the same $\sigma_\kappa$ and counts valid
transitions; reports $\phi$, the standardised $z$, and a binary
flag at $z>z_\alpha$ (\autoref{alg:detect-app},
\autoref{fig:refimpl}).

\paragraph{Token-level walkthrough.}
For the wiki prompt \texttt{Explain Donald Trump in a comprehensive
way.} on Llama-3.1-8B-Instruct at $S{=}5$, $\rho{=}0.5$, an
illustrative excerpt of the first $20$ generated tokens looks
roughly as below. We show the (token, state) pair at each position
and mark gated steps with $\bullet$ (the gate raised the mask;
ChainMark forced $s_{i+1} = (s_i{+}1)\bmod 5$) versus ungated steps
with $\circ$ (no mask; the model sampled freely):

\begin{quote}\small\ttfamily
\textbf{Donald}$_2$ $\circ$ \textbf{Trump}$_3$ $\circ$
(\textbf{born}$_4$ $\bullet$ \textbf{June}$_0$ $\bullet$
\textbf{14}$_1$ $\bullet$, \textbf{1946}$_2$ $\bullet$)
\textbf{is}$_3$ $\bullet$ \textbf{a}$_4$ $\circ$ \textbf{former}$_4$
$\circ$ \textbf{U.S.}$_0$ $\bullet$ \textbf{president}$_1$
$\bullet$ \textbf{and}$_2$ $\bullet$ \textbf{businessman}$_3$
$\bullet$ \dots
\end{quote}

At gated positions the next token's state is forced to be
$(\text{prev}+1)\bmod 5$ (here $S{=}5$, so the cycle is
$0{\to}1{\to}2{\to}3{\to}4{\to}0$); the model picks the highest-prob
\emph{token} whose state matches that requirement. At ungated
positions any token may follow. The gate threshold $\tau$ is
calibrated on a pilot so the realised gate rate tracks $\rho{=}0.5$.

\paragraph{Detection on this excerpt.}
A regulator running the detector
(\autoref{fig:refimpl}) on the same string with the same key $\kappa$
re-derives the state sequence
$2,3,4,0,1,2,3,4,4,4,0,1,2,3,\dots$ and
observes $11/14$ valid transitions, i.e.\ $\phi=0.79$. With
$p_0{=}0.20$ and $n{=}14$, $z = (0.79-0.20)/\sqrt{0.16/13} \approx
5.3$, well above $z_{0.01}=2.326$, so the detector flags
``watermarked''. A non-watermarked sample of the same length under
the null distribution averages $\phi\approx 0.20$ and $z\approx 0$.
